% !TEX encoding = UTF-8
% !TEX root = main.tex
\documentclass[a4paper,12pt]{extarticle}
\usepackage{geometry}
\usepackage[T1]{fontenc}
\usepackage[utf8]{inputenc}
\usepackage[english,russian]{babel}
\usepackage{mathptmx}
\usepackage{amsmath}
\usepackage{amsthm}
\usepackage{amssymb}
\usepackage{fancyhdr}
\usepackage{setspace}
\usepackage{graphicx}
\usepackage{caption}
\captionsetup{labelsep=endash, justification=raggedright, singlelinecheck=false}
\usepackage{colortbl}
\usepackage{tikz}
\usepackage{pgf}
\usepackage{subcaption}
\usepackage{listings}
\usepackage{indentfirst}
\usepackage[
backend=biber,
style=numeric,
maxbibnames=99
]{biblatex}
\addbibresource{refs.bib} % Убедитесь, что файл refs.bib находится в той же папке
\usepackage[colorlinks,citecolor=blue,linkcolor=blue,bookmarks=false,hypertexnames=true, urlcolor=blue]{hyperref}
\usepackage{indentfirst}
\usepackage{mathtools}
\usepackage{booktabs}
\usepackage[flushleft]{threeparttable}
\usepackage{tablefootnote}
\usepackage{chngcntr} % нумерация графиков и таблиц по секциям
\counterwithin{table}{section}
\counterwithin{figure}{section}

\makeatletter
% \renewcommand{@biblabel}[1]{1.} % Заменяем библиографию с квадратных скобок на точку:
\makeatother

\geometry{left=2.5cm}% левое поле
\geometry{right=1.0cm}% правое поле
\geometry{top=2.0cm}% верхнее поле
\geometry{bottom=2.0cm}% нижнее поле

\setlength{\parindent}{1.25cm}
\renewcommand{\baselinestretch}{1.5} % междустрочный интервал

\newcommand{\bibref}[3]{\hyperlink{1}{2 (3)}} % biblabel, authors, year
\addto\captionsrussian{\def\refname{Список литературы (или источников)}}

% Настройки нумерации списков
\renewcommand{\theenumi}{\arabic{enumi}}
\renewcommand{\labelenumi}{\arabic{enumi}}
\renewcommand{\theenumii}{.\arabic{enumii}}
\renewcommand{\labelenumii}{\arabic{enumi}.\arabic{enumii}.}
\renewcommand{\theenumiii}{.\arabic{enumiii}}
\renewcommand{\labelenumiii}{\arabic{enumi}.\arabic{enumii}.\arabic{enumiii}.}

% Путь к папке с изображениями (рекомендуется)
\graphicspath{{images/}}

\begin{document}

% --- Титульный лист ---
\begin{titlepage}
\begin{center}
    % Добавляем вертикальный отступ сверху, чтобы название не было прижато к краю
    \vspace*{6cm} 
    
    % НАЗВАНИЕ РАБОТЫ
    {\Huge\bfseries Прогнозирование за горизонтом прогнозирования: прогнозирование хаотических временных рядов\par}
    
    \vspace{4cm} % Отступ между названием и ключевыми словами
    
    % КЛЮЧЕВЫЕ СЛОВА
    \begin{minipage}{0.8\textwidth}
        \large
        \textbf{Ключевые слова:} Прогнозирование хаотических временных рядов, система Лоренца, расширение обучающей выборки, метод мотивов, непрогнозируемые точки, 
        кластеризация данных, дефицит данных.
    \end{minipage}

\end{center}

\vfill % Заполняет оставшееся пространство, отодвигая все, что ниже, к низу страницы
% Год и город убраны по требованиям анонимизации

\end{titlepage} % это титульный лист
\newpage

% --- Аннотация и ключевые слова ---
\input{sections/0_abstract.tex}
\newpage

% --- Оглавление ---
\tableofcontents
\newpage

% --- Основные разделы ---
\input{sections/1_introduction.tex}
\input{sections/2_literature_review.tex}
\input{sections/3_methodology.tex}
\input{sections/4_experiments.tex}
\input{sections/5_conclusion.tex}
\input{sections/6_acknowledgements.tex}

% --- Список литературы ---
\newpage
\printbibliography[title={Список использованных источников}]

\end{document}